% KYUSHU UNIVERSITY beamer template
% Created by Yu in 2024
% Modified for the report by the assistant.

\documentclass[11pt]{beamer}
% You can use \documentclass[11pt,aspectratio=169]{beamer}
% to adjust the aspect ratio to 16:9.

\usepackage{graphicx}
\usepackage{biblatex}


\usetheme{Madrid}

\author[Daniel Edwards Rodriguez]{Daniel Edwards Rodriguez}
\title[Graph-Aware SAT Solving]{Graph-Aware SAT Solving for Moderate-Scale Graph Colouring: A Performance-Robustness Trade-off Analysis}
\institute[KCL]{King's College London \\ Supervisor: Professor Hana Chockler}


\setbeamertemplate{navigation symbols}{}
\titlegraphic{\includegraphics[width=2cm]{kcl.png}}

\usepackage{Kyushu}

\begin{document}

\begin{frame}
\titlepage
\end{frame}

\begin{frame}{Introduction \& Motivation}
    \begin{itemize}
        \item \textbf{The Problem:} Graph colouring is a fundamental NP-complete problem with critical, real-world applications in scheduling, compiler design, and telecommunications.
        \item \textbf{The Gap:} Standard solvers are unreliable for mission-critical, moderate-scale problems (50-90 vertices). They often fail on structurally complex instances where reliability is essential.
        \item \textbf{Our Goal:} We shift the focus from average-case speed to \textbf{worst-case robustness}. We systematically trade a predictable performance overhead for a guarantee that a solution can be found.
    \end{itemize}
\end{frame}

\begin{frame}{The Problem: Baseline Solvers Are Brittle}
    \begin{itemize}
        \item Standard DPLL solvers fail unpredictably when faced with dense or highly-connected graphs.
        \item They suffer from an exponential explosion of the search space, which leads to timeouts.
        \item This makes them unsuitable for applications where failure is not an option.
    \end{itemize}
    \vfill
    \begin{block}{Our Hypothesis}
        Integrating domain-specific graph structural information can dramatically improve a SAT solver's robustness on these challenging problems.
    \end{block}
\end{frame}

\begin{frame}{Our Solution: A Robustness-Oriented Architecture}
    \begin{columns}[T]
        \begin{column}{0.5\textwidth}
            \begin{alertblock}{1. Graph-Aware Preprocessing}
                We invest a 15-20\% computational overhead to deeply analyze the graph's structure *before* solving begins.
            \end{alertblock}
            \pause
            \begin{block}{2. Centrality-Based Heuristics}
                We use metrics like degree and betweenness centrality to create a smart variable ordering that guides the solver effectively.
            \end{block}
        \end{column}
        \begin{column}{0.5\textwidth}
            \pause
            \begin{exampleblock}{3. Graceful Degradation}
                With built-in fallback mechanisms and circuit breakers, our solver is guaranteed to \textbf{never perform worse} than a baseline implementation.
            \end{exampleblock}
        \end{column}
    \end{columns}
    \vfill
    \uncover<4->{
    \centering
    \textit{This represents a strategic shift from an optimization-focused to a \textbf{reliability-focused} architecture.}
    }
\end{frame}

\begin{frame}{Key Results: Trading Speed for Reliability}
    \begin{columns}[T]
        \begin{column}{0.5\textwidth}
            \begin{block}{Predictable Overhead}
                Our solver has a consistent \textbf{1.45-1.47x} execution time overhead compared to the baseline. This cost is predictable and scales proportionally.
            \end{block}
        \end{column}
        \begin{column}{0.5\textwidth}
            \begin{alertblock}{100\% Reliability on Hard Problems}
                The enhanced solver achieves a \textbf{100\% success rate} on challenging instances where the baseline solver fails 25% of the time.
            \end{alertblock}
        \end{column}
    \end{columns}
    \vfill
    \begin{exampleblock}{A Qualitative Leap in Capability}
        On dense random graphs that cause the baseline solver to time out, our solver finds a solution in milliseconds. This is a fundamental difference in capability, not just a performance tweak.
    \end{exampleblock}
\end{frame}


\begin{frame}{Live Demo (3 minutes)}
	\frametitle{Live Demo: Robustness in Action}
	\begin{block}{Objective}
		To visually demonstrate the robustness of the enhanced solver on a graph instance specifically designed to make a standard solver fail.
	\end{block}
	
	\begin{columns}[T]
		\begin{column}{0.5\textwidth}
			\begin{alertblock}{Scenario 1: Baseline Solver}
				\begin{itemize}
					\item We will run a baseline DPLL solver on a dense, 30-vertex graph.
					\item \textbf{Expected Outcome:} The solver will quickly get stuck in the vast search space and time out.
				\end{itemize}
			\end{alertblock}
		\end{column}
		\begin{column}{0.5\textwidth}
			\begin{block}{Scenario 2: Our Enhanced Solver}
				\begin{itemize}
					\item We will run our graph-aware solver on the \textbf{exact same problem}.
					\item \textbf{Expected Outcome:} The solver will find a solution almost instantly.
				\end{itemize}
			\end{block}
		\end{column}
	\end{columns}
	\vfill
	\begin{center}
		\textit{This demo highlights the practical value of our approach: trading a small, consistent overhead for a 100\% success guarantee on difficult problems.}
	\end{center}
\end{frame}


\begin{frame}{Conclusion \& Future Work}
    \begin{block}{Key Contribution}
        We have validated that graph-aware SAT solving is a viable and powerful strategy for reliability-critical applications. A predictable performance trade-off can yield enormous robustness gains.
    \end{block}
    \begin{columns}[T]
        \begin{column}{0.5\textwidth}
            \begin{exampleblock}{\makebox[0.5\textwidth][c]{Practical Impact}}
                \begin{itemize}
                    \item Academic Scheduling
                    \item Compiler Design for Safety-Critical Systems
                    \item Telecommunications
                \end{itemize}
            \end{exampleblock}
            \vspace{0.5cm}
        \end{column}
        \begin{column}{0.5\textwidth}
            \begin{block}{\makebox[0.5\textwidth][c]{Future Work}}
                \begin{itemize}
                    \item Develop adaptive strategies.
                    \item Explore hybrid approaches.
                    \item Extend evaluation to larger graphs.
                \end{itemize}
            \end{block}
            \vspace{0.5cm}
        \end{column}
    \end{columns}
\end{frame}

\end{document}